\section{System Description} Havven is designed to incentivise stability in a decentralised cryptocurrency denominated in some external currency, such as the USD. The dual-token system is combined with a set of novel incentive mechanisms that stabilise the price of one of the two tokens. The first token is the havven token - with a lowercase h - to avoid confusion with the Havven system itself. The second token is the nomin. \\

\noindent The havven token serves two functions:

\begin{enumerate}
\item{To provide the system with collateral (the system itself is tokenised).}
\item{To allow actors to contribute to the price stabilisation process through a set of incentives.}
\end{enumerate}

\noindent The stablecoin is the nomin. The purpose of the nomin is to track the price of a chosen external currency via the actions of havven token holders; holders of havven participate in modifying the level of supply in the nomin market such that the market price of nomin is maximally stable.

\subsection{Incentive Layering}

We classify the various incentives that can be applied in a stablecoin system. Note that any subset of these can be linearly combined in order to produce a sophisticated and powerful incentive structure. Havven's approach to achieving price stability is to be as passive as possible and only switch on higher levels of incentivisation when necessary. The order in which these categories appear is the order in which they are applied:

\paragraph{Overcollateralisation}

The basis for price stability within Havven is overcollateralisation of the stablecoin. This means the value of escrowed havven must be strictly greater than the value of all nomin in circulation. In Havven, this ratio of nomin value to havven value is known as the utilisation ratio.

\paragraph{Fees}

\noindent The second layer of economic incentives is to provide havven holders with fees in accordance with their performance in adjusting the nomin supply. These fees are generated from small charges on all nomin transactions. The fees are directed to the havven holders as a reward for helping maintain the correct nomin supply.

\paragraph{Interest Rates}

\noindent Interest rates on escrowed havven can be applied in addition to the application of fees, in either fixed or floating havven supply regimes. Interest rates will be discussed in a future iteration of the whitepaper.

\paragraph{Collateral Recovery}

\noindent As a final layer of incentives, forced recovery of an actor's escrowed havven may be required in order to equilibriate individual positions of utilisation ratio. Collateral recovery will be discussed in a future iteration of the whitepaper.

\subsection{Overcollateralisation}

\noindent We first introduce the core system variables:

\begin{align*}
H &= \text{Quantity of havven,} & N &= \text{Quantity of nomin,} \\
P_h &= \text{havven price,}  & P_n &= \text{nomin price.}
\end{align*}

\noindent All havven tokens are created in the initial system state, so $H$ is constant. The quantity of nomin, $N$, floats in response to the actions of havven holders. The Havven system needs to incentivise havven holders to maintain $N$ such that the nomin price, $P_n$, is stable at \$1.

\subsubsection{Nomin Issuance}

\noindent Nomins can only be issued when a havven holder decides to escrow some or all of their havven. In Havven, escrowing means locking some quantity of havvens $ H_e $ in order to generate a quantity of nomins $q$, such that:
\begin{align*}
q &= P_h * H_e * U_{max}, \\
& U_{max} < 1. \\
\end{align*}

\noindent In the above equation, $U_{max}$ represents what is known as the maximum utilisation ratio, which will be discussed below. By setting $U_{max} < 1 $, all nomins are overcollateralised when issued. \\

\noindent After generating the nomins, the system places a \textbf{limit sell} order with a price of \$1 on a decentralised exchange. This means that the nomin will be sold at the current market price, down to a minimum price of \$1 USD. If we assume implementation on Ethereum, then the nomin are sold for ETH, with the proceeds of the sale remitted to the issuer. \\

\subsubsection{Releasing Havvens from Escrow}

\noindent In order to access the original havven that have been escrowed, the issuer must return the same quantity of nomin to the system for destruction. This is known as \textbf{burning} the nomins. Nomin can be purchased in the open market.

\newpage

\subsubsection{Issuance Example (Diagrams to come)}

\noindent 1. Bob purchases 10 havven at \$10 each, total value \$100. We can visualise Bob's wallet below.

\begin{center}
\begin{tikzpicture}
\fill[color=black] (0,0) rectangle (10,1.0);
\end{tikzpicture}
\end{center}

\noindent 2. Bob decides to issue some Nomins. He escrows 5 of his havvens worth \$50, locking them in the system. \\

\begin{center}
\begin{tikzpicture}
\fill[color=black] (0,0) rectangle (10,1.0);
\end{tikzpicture}
\end{center}

\noindent 3. The maximum utilisation ratio is 0.2. The system generates and sells 10 nomin on the market and transfers the proceeds, in ETH, to Bob's wallet. \\

\begin{center}
\begin{tikzpicture}
\fill[color=black] (0,0) rectangle (10,1.0);
\end{tikzpicture}
\end{center}

\noindent 4. In order to release the escrowed havven, Bob must return 10 nomin.

\begin{center}
\begin{tikzpicture}
\fill[color=black] (0,0) rectangle (10,1.0);
\end{tikzpicture}
\end{center}

\subsubsection{Price changes}

\noindent \emph{What happens if the price of havven changes?} \\

\noindent All issuance of nomin is done at the current $P_h$. However, when $P_h$ changes, the quantity of escrowed havvens changes with it (not the \emph{value}). An increase in $P_h$ means that fewer of Bob's havvens are escrowed. By contrast, a decrease in the $P_h$ means that more of his havvens are escrowed. This process occurs automatically in order to ensure that the system remains overcollateralised. \\ 

\noindent \emph{What happens if the price of nomin changes?} \\ 

\noindent In order to release escrowed havven, Bob must return the same quantity of nomin that he issued. This means that if $P_n$ has increased in the market, he will need to spend more ether than he received when he issued in order to release his havven. Conversely, if $P_n$ has decreased, Bob will need to spend less in order to release his havvens.

\newpage

\subsubsection{Global Utilisation Ratio}

\noindent Havven can ensure that nomins are overcollateralised at the time of issuance, however, a reduction in $P_h$ lowers the value of the escrowed havvens. The global utilisation ratio is defined as the total value of nomins divided by the total value of havvens:

$$ U = \frac{P_n * N}{P_h * H}.$$ \\

\noindent Intuitively, if $U = 1$, the value of nomins and havvens are equal. Hence, given our overcollateralisation property, the system targets $U <  1$. To do this, Havven only allows the issuance of nomin up to a maximum utilisation ratio.

\subsubsection{Collateralisation Target}

\noindent Given the need to adjust the supply of nomins, an optimum utilisation ratio is defined as the point at which maximum incentives are applied:

$$ 0 \leq U_{opt} \leq U_{max} \leq 1.$$

\noindent Each havven holder that has issued nomins has a unique utilisation ratio, $ U_i $. The system can measure the degree to which their $ U_i $ is above or below the optimum and reward them accordingly. In this way the system incentivises the issuance and destruction of nomins. $ U_{opt} $ is defined formally below in terms of $ P_n $. As  the nomin price diverges from the desired \$1, the incentive to either expand or contract nomin supply increases.

\newpage
\subsection{Transaction Fees} In addition to overcollateralisation, the system needs a direct incentive mechanism that can react to changes in the nomin price. 

\subsubsection{Nomin Equilibrium} Demand and supply economics shows that there exists some optimal nomin supply where the related level of demand yields an equilibrium price of \$1. This quantity can be expressed in terms of a special utilisation ratio, $U_{opt}$. The graph below visualises this situation. \\

\begin{center}
\begin{tikzpicture}[scale=3]

% draw axes
\draw [<->, thick] (0,2) node (yaxis) [above] {$P$} |- (2.5,0) node (xaxis) [right] {$Q$};

% draw intersecting lines
\draw (0.5, 0.5) coordinate (a_1) -- (2,1.8) coordinate (a_2);
\draw (0.5, 1.8) coordinate (b_1) -- (2,0.5) coordinate (b_2);

% calculate coordinate of intersection
\coordinate (c) at (intersection of a_1--a_2 and b_1--b_2);

\draw[dashed] (yaxis |- c) node [left] {$1$} -| (xaxis -| c) node[below] {$u_{opt}$};

\end{tikzpicture}
\end{center}
	
\subsubsection*{Demand}

\noindent The system is unable to influence the demand for nomin. We assume that some level of demand exists given the utility of nomin as a stable cryptocurrency.

\subsubsection*{Supply}

\noindent However where demand is unable to be directly influenced, the supply of nomin is controlled by havven holders who use the system to issue and burn nomin in response to its incentives. Maximum incentive is achieved when $U_i = U_{opt}$, such that $P_n = 1$. $U_{opt}$ will be discussed further below. \\

\newpage
\subsubsection{Nomin transaction fees} Every time a nomin transaction occurs, the Havven system charges a small transaction fee. Transaction fees allow the system to generate revenue, which it can distribute to havven holders as an incentive to maintain nomin supply at $U_{opt}$. \\

\noindent The fee rate charged on nomin transactions is $\alpha_c$. It is constant and will be sufficiently small that it provides little to no friction for the user.\\

$$ a_c = k.$$ 

\begin{center}
\begin{tikzpicture}[scale=3]

% draw axes
\draw [<->, thick] (0,2) node (yaxis) [above] {$\alpha_c$} |- (2.5,0) node (xaxis) [right] {$P_N$};

% draw one line
\draw (0.0, 0.75) coordinate (a_1) -- (2.0,0.75) coordinate (a_2);

\coordinate (c) at (1.75, 0.75);

\draw[dashed] (yaxis |- c) node [left] {k};

\end{tikzpicture}
\end{center}

\newpage
\subsubsection{Fees received by Havven Holders}

\noindent The fee rate paid to a havven holder that has escrowed is $\alpha_r$. This rate changes with respect the individual's unique utilisation ratio, $U_i$. It increases linearly to a maximum at the optimal utilisation ratio $U_{opt}$, before quickly diminishing as $U_i$ approaches $U_{max}$. Beyond the maximum utilisation ratio $\alpha_r$ is 0. Note, $\alpha_r$ is applied to the pool of collected fees which is determined by $\alpha_c$. \\

$$ \alpha_{R,t,I} = \alpha_{base,t} * f_{i,t}(U_{i,t}, U_{opt}, U_{max,t}), $$

\[
f_{i,t}(U_{i,t}, U_{opt}, U_{max,t}) = 
\begin{cases}
 \frac{U_{i,t}}{U_{opt,t}} &\mbox{when } U_i \leq U_{opt}, \\[1em]
 \frac{U_{max,t} - U_{i,t}}{U_{max,t} - U_{opt,t}} &\mbox{when } U_{opt} \leq U_i \leq U_{max}, \\[1em]
 0 &\mbox{otherwise}.
 \end{cases}
\]

\begin{center}
\begin{tikzpicture}[scale=3]

% draw axes
\draw [<->, thick] (0,2) node (yaxis) [above] {$\alpha_r$} |- (2.5,0) node (xaxis) [right] {$U_i$};

% draw two lines
\draw (0.0, 0.0) coordinate (a_1) -- (1.75,0.75) coordinate (a_2);
\draw (1.75, 0.75) coordinate (a_1) -- (2.25,0.0) coordinate (a_2);

\coordinate (c) at (1.75, 0.75);

\draw[dashed] (yaxis |- c) node [left] {$a_{base}$} -| (xaxis -| c) node[below] {$u_{opt}$};

\end{tikzpicture}
\end{center}

\noindent This fee distribution curve encourages havven holders who have escrowed to maintain their $U_i$ at $U_{opt}$.  \\

\newpage

\subsubsection{Deriving the base fee rate} The total amount of fees collected from users has to be equal to the total amount of fees paid to the havven holders. We define the total amount of fees collected, $A_{c,t}$ below: \\
$$ A_{c,t}  = \alpha_{c,t} * v_{n,t} * \sum\limits_I N_{I,t}, $$
$$ v_{n,t} \text{ the velocity of nomins at t}, $$
$$  \sum\limits_I N_{I,t} \text{ the total issued nomins}. $$
Next we define the total fees paid to havven holders $A_{R,t}$: \\
$$ A_{R,t} = \sum\limits_I \alpha_{R,t,I}. $$
Havven requires that: \\
$$ A_{R,T} =  A_{c,t}. $$ \\
Substituting our earlier definition of $\alpha_{R,t,I}$: \\
$$ \alpha_{base,t} *\sum\limits_I f_{i,t}(U_{i,t}, U_{opt}, U_{max,t}) =  \alpha_{c,t} * v_{n,t} * \sum\limits_I N_{I,t}. $$
Solving for $\alpha_{base,t}$:\\
$$ \alpha_{base,t} = \frac{\alpha_{c,t} * v_{n,t} * \sum\limits_I N_{I,t}}{\sum\limits_I f_{i,t}(U_{i,t}, U_{opt}, U_{max,t})}.$$

\noindent We have now defined the maximum fee rate, $\alpha_{base}$, in terms of the fees collected, $A_{c,t}$. This rate should be achieved when an individuals $U_i$ is at $U_{opt}$. \\

\noindent The definition of $U_{opt}$ must therefore provide the following incentive. If $P_n > 1$ then the system must encourage more nomin to be issued. If $P_n < 1$, the system must encourage nomin to be burned. 

\newpage

\subsection{Utilisation Ratio}
\subsubsection{Optimal Utilisation Ratio}

\noindent The optimal utilisation ratio $U_{opt}$ is a target for havven holders to reach in order to maximise the amount of fees they receive. $U_{opt}$ is defined in terms of $P_n$ such that havven holders can influence the price of nomin through directly controlling the supply of nomin (a havven holder can change their individual utilisation ratio by buying or issuing more nomins). \\
 
\noindent The function for $U_{opt}$ given below provides our dynamic target for havven holders based on the price of nomin. The curve shows that the when $P_n$ is close to \$1, $ f'(P_n) $ is small. However, the further $P_n$ diverges from \$1, the larger the derivative becomes, providing an increasing incentive (via fees) for a havven holder to move toward $U_{opt}$.

$$ U_{opt} = f(P_n) * U,$$
$$ f(P_N) = max(\sigma * (x - 1)^{\phi} + 1, 0), $$
$$\text{where } 0 \leq \sigma, \text{ the price sensitivity parameter}, $$
$$\phi \geq 1, \text{ the flattening parameter}. $$ \\

\begin{center}
\begin{tikzpicture}[scale=1.15]
\begin{axis}[
    axis lines = left,
    xlabel = $P_N$,
    ylabel = {$f(P_N)$},
    xtick = {1},
    ytick = {0},
    y label style={at={(axis description cs:.15,1.00)}, rotate=270,anchor=south},
    x label style={at={(axis description cs:1.05,.13)},anchor=north},
]
\addplot[
domain=0:2,
range=0:2
]
{(x-1)^3 + 1};
\addplot [
domain=0:2,
range=0:2,
dashed
]  
 {x};
 \addplot[dashed, samples=50, smooth,domain=0:6] coordinates {(1,1)(1,0)};
\end{axis}
\end{tikzpicture}
\end{center}

\newpage

\subsubsection{Maximum Utilisation Ratio}

\noindent Havven seeks to maintain $U < U_{max} < 1$, in order to remain overcollateralised. It might seem intuitive that $U_{max}$ should be a static value. However, since $U_{opt}$ changes linearly with $P_n$ and inversely with $P_h$, there are several situations where $U_{max}$ may need to change. Below we define $U_{max}$. \\

\[
U_{max} = 
\begin{cases}
 U_{base} &\mbox{when } U_{opt} \leq U_{base}, \\[1em]
 a * U_{opt} &\mbox{otherwise}.
 \end{cases}
\]

\begin{center}
\begin{tikzpicture}[scale=3]

% draw axes
\draw [<->, thick] (0,2) node (yaxis) [above] {$U_{max}$} |- (2.5,0) node (xaxis) [right] {$U_{opt}$};

% draw intersecting lines
\draw (0.0, 0.75) coordinate (a_1) -- (0.75,0.75) coordinate (a_2);
\draw (0.75, 0.75) coordinate (a_1) -- (1.8,1.8) coordinate (a_2);

% calculate coordinate of intersection
\coordinate (c) at (0.75,0.75);

\draw[dashed] (yaxis |- c) node [left] {$U_{base}$} -| (xaxis -| c) node[below] {$U_{base}$};

\end{tikzpicture}
\end{center}

\newpage

\subsection{Intrinsic Havven Price} With the havven token being ERC20 compliant, it will have a market price on both decentralised and centralised exchanges. \\

\noindent While the Havven system will access the current market price via a price oracle, it is beneficial to define a $P_h$ that can be determined internally to avoiding the influence of speculation. Ignoring speculative demand, $P_h$ can be expressed as a function of the transaction fees that the system charges. Below we define an initial iteration of the intrinsic $P_h$.

\begin{align*} 
P_{h,t} &= \frac{1}{H}* \sum\limits_{t=1}^\infty \frac{d_{n,t} *v_{n,t} * \alpha_{R,t}}{(1+R)^t} \approx \frac{d_{n,t} *v_{n,t} * \alpha_{R,t}}{R * H}, \\
& P_{h,t} \text{ is the price of one havven at time } t, \\
& H \text{ is the number of havvens}, \\
& d_{n,t} \text{ is the demand for nomins at t}, \\
& v_{n,t} \text{ is the velocity of nomins at t}, \\
& \alpha_{R,t} \text{ is the fee from trade with nomins}, \\
& R \text{ is the interest rate / rate of return of havvens}. \\
\end{align*}

\newpage